\documentclass[aspectratio=169, UTF8]{ctexbeamer}
\makeatletter
\def\input@path{{./}{../}{theme/}{../theme/}}
\makeatother
\usetheme{CX}

\title{CXSlide 最小示例}
\subtitle{验证骨架可工作}
\author{作者}
\institute{某某大学}
\date{\today}

\begin{document}

\begin{frame}[plain, noframenumbering]
  \titlepage
\end{frame}

\begin{frame}{目录}
  \tableofcontents
\end{frame}

\section{引言}

\begin{frame}{列表与强调}
  \begin{itemize}
    \item 第一点
    \item \alert{强调}内容
    \item 嵌套列表
      \begin{itemize}
        \item 子项 A
        \item 子项 B
      \end{itemize}
  \end{itemize}
\end{frame}

\section{方法}

\begin{frame}{区块}
  \begin{block}{定理}
    $|\langle \mathbf{u}, \mathbf{v} \rangle|^2
      \leq \langle \mathbf{u}, \mathbf{u} \rangle
           \cdot \langle \mathbf{v}, \mathbf{v} \rangle$
  \end{block}
  \begin{alertblock}{注意}
    等号条件：线性相关。
  \end{alertblock}
  \begin{exampleblock}{例}
    $\mathbf{u}=(1,2),\;\mathbf{v}=(3,4)$。
  \end{exampleblock}
\end{frame}

\begin{frame}{公式}
  \begin{equation}
    e^{i\pi} + 1 = 0
  \end{equation}
  \begin{equation}
    \hat{f}(\xi) = \int_{-\infty}^{\infty}
      f(x)\, e^{-2\pi i x \xi}\, dx
  \end{equation}
\end{frame}

\end{document}
